\section{From Monitoring to Observability}

A traditional monitoring system usually starts with simple questions. Is the service up? Is the
disk full? Is the CPU high? These questions are necessary, but they are not enough for a facility
like Rubin. Observability enables understanding of a system's internal behavior from the signals it
emits. In Kubernetes, observability is typically built from metrics, logs, and traces, along with
the tools that make them usable for operators \cite{kubernetes-observability}.

The practical difference is important. Monitoring can tell us that an alert pipeline is slow.
Observability should help us decide whether the reason is the Kafka broker, the network, the
storage layer, an overloaded Kubernetes node, an ingress problem, a database, or something related
to the observing sequence. Monitoring can indicate that a pod has restarted; observability should
help us understand what changed around that time, who was affected, and what to check next.

Our approach is also a journey. The current platform is not the first iteration we tried, and it
will not be the last. The important part is that the platform became more reproducible and easier
to inspect over time.

At the beginning, we did what most teams do: we installed Icinga and set it up for basic up/down
checks. It was a reasonable first step, but it did not take long to realize that a platform of this
complexity could not be understood solely through availability probes. We could see things were
failing, but we had no insight into why.

Among the team, we had backgrounds with Nagios, Zabbix, Zenoss, and even HP OpenView, and more,
but none of those felt like the right direction. We decided that open-source tools were the right
choice, not because they are free of cost, but because they can be operated, modified, and
connected to the rest of the observatory without vendor lock-in.

We knew we needed outside help to build a modern observability stack, so we hired a contractor.
One of their consultants was an active maintainer of Prometheus, which gave us confidence that we
were talking to the right people.

We wrote down our requirements and handed them over. What we learned over the following months is
that our environment is genuinely unusual. Our components come from many different vendors and do
not follow any standard pattern. That made everything slower and harder to implement than either
side expected. It took close to two years before the contractor delivered something we could call
substantial. Even then, we were never happy with the logging side. They had gone with OpenSearch,
and while we understood the reasoning, it never felt like a good fit for us. It was not well
integrated with Grafana, and the resource consumption was hard to justify. The dashboards were also
a source of friction. We eventually came to understand that part of the problem was our own
requirements; we had asked for dashboards built around a preconceived idea of what should be shown,
when in reality dashboards need to be driven by what is actually happening in the infrastructure,
not the other way around.

When the contract ended, we took the chance to redo the logging from scratch. We picked Loki,
mostly because it is part of the Grafana Labs ecosystem and fits naturally with everything else we
were already running. That choice felt like the moment we stopped working around someone else's
observability architecture and started building something that actually made sense for us.